\scanpage{9}\mspage{8}
$\inf$ dans $\unsure{\mathcal{F}|\mathfrak{F}}$ est $\varnothing_{\mathcal{F}}$, le plus petit élément. de $\mathcal{F}$ ou \uline{figure vide}. Ainsi
\[
F\text{ et }G\text{ disjoints}\quad\Longrightarrow\quad F\text{ et }G\text{ compatibles}.
\]

On a, par les axiomes:
\[
\begin{aligned}
&&\text{\deleted{tout}}\quad\inserted{\operatorname{omb}(F)\text{ et }\operatorname{omb}(G)}\\
F\text{ et }G\text{ disjoints}&\quad\Longleftrightarrow\quad&\deleted{\operatorname{omb}(F)\cap\operatorname{omb}(G)}\\
&&\text{sont ``disjoints'' au sens}\\
&&\text{suivant: ils sont \unsure{disjoint|disjoints}}\\
&&\text{ensemblistement i.e.}\\
&&\operatorname{omb}(F)\cap\operatorname{om}(G)=\varnothing,\ \text{et}\\
&&\text{tout élément de }\operatorname{omb}(F)\text{ est}\\
&&\text{compatible avec tout elt}\\
&&\text{de }\operatorname{om}(G)
\end{aligned}\tag*{(1.13)}
\]
--- i.e. deux figures sont disjointes ssi elles n'ont pas de lieu communs, et tout lieu de l'une est \deleted{disjoint} \inserted{\deleted{com} disjoint} (au sens de l\deleted{a}\inserted{'}\deleted{théorie}atelier) de tout lieu de l'autre;

Cela montre que les relations de \deleted{dis}compatibilité et de disjonction sont importantes aussi dans les \deleted{\unsure{en des|en de}} \inserted{sous-}ensembles $\mathcal{M},\mathcal{L}$ de $\mathcal{F}$. Dans le cas des ateli\unsure{ers} ensemblistes, ces relations sont ``triviales'' sur $\mathcal{L}$: deux lieux sont tjrs \deleted{disjoints} \inserted{compatibles}, et disjoints ssi ils sont distincts. Par contre, même dans le cas ensembliste et $\overline m$ modéré, la relation
% The ending of ateliers is retraced; a separate canceled older ending cannot be securely isolated.
